\scanpage{124}
\mspage{39}
% archive sheet number: 123

% L01
sur $X_0$, et à l'opération de $G|X_0$ sur ces faisceaux.

% L02-07
Pour toute composante irréductible $D_i$ de $D$, soit $m_i$ la multiplicité de \deleted{\unsure{ce}} $D_i$ dans $\operatorname{div}(\pi)$; Soit $m$ le ppcm des $m_i$, et soit $m'=m$ si $X$ de car. nulle, $m'=$ partie de $m$ première à $p=\operatorname{car}V(\pi)$ dans le cas contraire. Écrivons $G$ additivement et considérons le groupe proétale

% L08
\[
\deleted{\unsure{mG}}\qquad G'=m'G\inserted{=mG}\subset G
\]

% L09
sur $X$. \deleted{\unsure{qui}}

% L10-16
\uline{Corollaire} 4.17. Dans les \deleted{hypothèses}\inserted{hypothèses et notations} de 4.16, supposons de plus que pour tt localisé strict $\overline{X}^{\unsure{s|0|6}}$ de $X$ en un point de $X_0$, les $\overline{X}(m)$ définis en termes du diviseur induit $\overline{D}$ (cf. 4.8) satisfont au th. de pureté. Alors l'opération de $G'|X_0$ sur \deleted{\unsure{les}} $R^*f_*(\Lambda_{U'})|X_0$ est triviale.

% L17-18
Cela résulte aussitôt de 4.15 par localisation stricte.

% L19-21
% insertion attachment uncertain: phrase is between structure de and the R^*f_* line
\deleted{4.18. En fait, on peut préciser la structure de \inserted{comme faisceau d'anneaux gradués à groupe d'op. $G|X_0$} $R^*f_*(\Lambda_{U'})|X_0$, On trouve tout d'abord}

% L22
\[
\deleted{R^*f}
\tag*{\deleted{(4.18.1)}}
\]
